% ---------------------------------------------------------------------
%  CHƯƠNG 6 — KẾT LUẬN VÀ HƯỚNG PHÁT TRIỂN
% ---------------------------------------------------------------------
\chapter{Kết luận và hướng phát triển}
\label{ch:ketluan}

\begin{muctieu}
Tổng kết toàn bộ quá trình thực hiện bài tập lớn, đúc kết các bài học kỹ thuật
quý giá và đề xuất các hướng mở rộng phát triển hệ thống trong tương lai.
\end{muctieu}

\section{Kết luận}

Sau quá trình nỗ lực nghiên cứu và triển khai, nhóm sinh viên gồm \textbf{Trần
Thanh Bình} và \textbf{Huỳnh Tiến Đạt} đã hoàn thành xuất sắc bài tập lớn môn
Lập trình Java nâng cao (AJP20101) với Đề tài số 30: ``Trò chơi đoán số nhiều
người qua mạng TCP''. Hệ thống được thiết kế theo đúng các chuẩn mực kỹ thuật
phần mềm hiện đại, có tính mở rộng cao và vận hành ổn định trong môi trường
mạng đa luồng.

\section{Bài học rút ra}

Qua quá trình thực hiện bài tập lớn, nhóm đã rút ra được nhiều bài học kinh
nghiệm sâu sắc:
\begin{enumerate}[leftmargin=1.6em, itemsep=2pt]
  \item \textbf{Sức mạnh của phương pháp TDD:} Việc viết test trước giúp định
        hình rõ ràng giao diện hàm và lường trước các trường hợp biên, tiết
        kiệm đến 70\% thời gian gỡ lỗi so với cách viết mã thông thường.
  \item \textbf{Giá trị của kiến trúc phân tầng sạch (Clean Architecture):}
        Tách rời tầng nghiệp vụ khỏi socket mạng giúp kiểm thử logic độc lập,
        dễ dàng thay đổi giao thức mạng mà không ảnh hưởng tới luật chơi.
  \item \textbf{Kỹ năng kiểm soát luồng và Deadlock:} Hiểu rõ bản chất của
        điều phối luồng trong hệ điều hành, vận dụng đúng công cụ đồng bộ hóa
        (khóa đối tượng hẹp, Lock Ordering) giúp ứng dụng luôn an toàn và ổn
        định.
\end{enumerate}

\section{Hướng phát triển}

Một số hướng phát triển mở rộng trong tương lai:
\begin{itemize}[leftmargin=1.6em, itemsep=2pt]
  \item Tích hợp cơ sở dữ liệu quan hệ (PostgreSQL / SQLite) để lưu trữ vĩnh
        viễn bảng điểm và lịch sử người chơi qua nhiều mùa giải.
  \item Phát triển giao diện Web tương thích WebSocket để người chơi có thể
        tham gia trực tiếp từ trình duyệt điện thoại hoặc máy tính mà không
        cần cài đặt ứng dụng Swing.
  \item Ứng dụng mô hình bảo mật SSL/TLS (Java \code{SSLSocket}) để mã hóa
        toàn bộ dữ liệu trao đổi giữa máy khách và máy chủ.
\end{itemize}
